\documentclass[11pt,a4paper]{article}
\usepackage[utf8]{inputenc}
\usepackage[italian]{babel}
\usepackage{amsmath}
\usepackage{graphicx}
\usepackage{listings}
\usepackage{xcolor}
\usepackage{hyperref}
\usepackage{geometry}
\usepackage{enumitem}
\usepackage{booktabs}
\usepackage{longtable}
\usepackage{titlesec}

\geometry{margin=2.5cm}

% Configurazione colori per codice
\definecolor{codegreen}{rgb}{0,0.6,0}
\definecolor{codegray}{rgb}{0.5,0.5,0.5}
\definecolor{codepurple}{rgb}{0.58,0,0.82}
\definecolor{backcolour}{rgb}{0.95,0.95,0.92}

\lstdefinestyle{mystyle}{
    backgroundcolor=\color{backcolour},
    commentstyle=\color{codegreen},
    keywordstyle=\color{magenta},
    numberstyle=\tiny\color{codegray},
    stringstyle=\color{codepurple},
    basicstyle=\ttfamily\footnotesize,
    breakatwhitespace=false,
    breaklines=true,
    captionpos=b,
    keepspaces=true,
    numbers=left,
    numbersep=5pt,
    showspaces=false,
    showstringspaces=false,
    showtabs=false,
    tabsize=2,
    frame=single,
    frameround=tttt
}

\lstset{style=mystyle}

% Formattazione sezioni
\titleformat{\section}
{\Large\bfseries}
{}
{0em}
{}[\titlerule]

\titleformat{\subsection}
{\large\bfseries}
{}
{0em}
{}

% Informazioni documento
\title{
    \textbf{Documentazione Tecnica: Sistema UserActivity} \\
    \large{Analisi Completa del Componente di Tracking e Gamification}
}
\author{ExtraTracker Development Team}
\date{\today}

\begin{document}

\maketitle

\begin{abstract}
Questo documento fornisce una documentazione tecnica completa ed esaustiva del sistema \texttt{UserActivity}, componente centrale per il tracking delle attività utente e l'implementazione della gamification in ExtraTracker. Il documento analizza l'architettura, i punti di forza, le integrazioni, i limiti attuali e le aree di miglioramento, fornendo una visione professionale e dettagliata del sistema.
\end{abstract}

\tableofcontents
\newpage

\section{Introduzione}

\subsection{Panoramica}

Il sistema \texttt{UserActivity} è il cuore del meccanismo di tracking e gamification di ExtraTracker. Esso registra ogni interazione significativa dell'utente con la piattaforma, trasformandola in dati strutturati utilizzabili per:

\begin{itemize}
    \item \textbf{Gamification}: Calcolo di XP, livelli e streak
    \item \textbf{Analytics}: Analisi comportamentale e pattern di utilizzo
    \item \textbf{Insights}: Generazione di consigli personalizzati (AI Coach)
    \item \textbf{Dashboard}: Visualizzazione di attività recenti e statistiche
\end{itemize}

\subsection{Scopo del Documento}

Questo documento mira a:

\begin{enumerate}
    \item Fornire una comprensione completa dell'architettura del sistema
    \item Documentare tutti i componenti e le loro interazioni
    \item Identificare punti di forza e vantaggi architetturali
    \item Evidenziare limiti, bug e aree di miglioramento
    \item Fornire linee guida per manutenzione e sviluppo futuro
\end{enumerate}

\section{Architettura del Sistema}

\subsection{Componenti Principali}

Il sistema UserActivity è composto da quattro componenti principali:

\begin{enumerate}
    \item \textbf{UserActivity Model}: Modello MongoDB per la persistenza
    \item \textbf{ActivityService}: Servizio centrale per la logica di business
    \item \textbf{ActivitySubscriber}: Subscriber per il pattern Observer
    \item \textbf{ActivityQueue}: Coda asincrona con retry automatico
\end{enumerate}

\subsection{Diagramma Architetturale}

Il flusso di registrazione attività segue questo pattern:

\begin{verbatim}
┌─────────────────┐
│  Service Layer  │  (goalService, studyService, etc.)
│  (Business)     │
└────────┬────────┘
         │ emit('event')
         ▼
┌─────────────────┐
│   Event Bus     │  (Pattern Observer)
│  (Decoupling)   │
└────────┬────────┘
         │ on('event')
         ▼
┌─────────────────┐
│ActivitySubscriber│  (Converte eventi in attività)
└────────┬────────┘
         │ addActivityJob()
         ▼
┌─────────────────┐
│ ActivityQueue   │  (Bull Queue con retry)
└────────┬────────┘
         │ process()
         ▼
┌─────────────────┐
│ActivityService  │  (Logica gamification)
└────────┬────────┘
         │ create()
         ▼
┌─────────────────┐
│ UserActivity    │  (MongoDB)
│     Model       │
└─────────────────┘
\end{verbatim}

\section{UserActivity Model}

\subsection{Schema MongoDB}

Il modello \texttt{UserActivity} definisce la struttura dei dati per ogni attività registrata:

\begin{lstlisting}[language=JavaScript, caption=Schema UserActivity]
const userActivitySchema = new mongoose.Schema({
    type: {
        type: String,
        enum: [
            'SESSION_COMPLETE',
            'GOAL_CREATED',
            'GOAL_COMPLETED',
            'GOAL_ARCHIVED',
            'WORK_SESSION_LOGGED',
            'WORK_NOTE_CREATED',
            'HABIT_CHECKIN',
            'HABIT_SKIPPED',
            'PROJECT_CREATED',
            'PROJECT_COMPLETED',
            'STREAK_BONUS',
        ],
        required: true,
    },
    entityId: {
        type: mongoose.Schema.Types.ObjectId,
        default: null,
        index: true,
    },
    category: {
        type: String,
        trim: true,
        default: '',
        index: true,
    },
    sentiment: {
        type: Number,
        min: 1,
        max: 5,
        default: null,
    },
    xpEarned: {
        type: Number,
        default: 0,
        min: 0,
    },
    metadata: {
        type: mongoose.Schema.Types.Mixed,
        default: {},
    },
}, {
    timestamps: { createdAt: true, updatedAt: false },
});
\end{lstlisting}

\subsection{Indici}

Il modello utilizza i seguenti indici per ottimizzare le query:

\begin{itemize}
    \item \texttt{\{ user: 1, createdAt: -1 \}}: Query per attività recenti di un utente
    \item \texttt{\{ user: 1, type: 1, createdAt: -1 \}}: Query filtrate per tipo
    \item \texttt{\{ user: 1, category: 1, createdAt: -1 \}}: Query filtrate per categoria
    \item \texttt{\{ 'metadata.entityId': 1, createdAt: -1 \}}: Query per entità specifica
\end{itemize}

\textbf{Punto di Forza}: Gli indici sono ben progettati per le query più comuni (dashboard, analytics).

\subsection{Multi-Tenancy}

Il modello utilizza il plugin \texttt{multiTenancyPlugin} per garantire isolamento dei dati tra utenti:

\begin{lstlisting}[language=JavaScript]
userActivitySchema.plugin(multiTenancyPlugin, {
    tenantField: 'user',
    required: true,
});
\end{lstlisting}

\textbf{Punto di Forza}: Isolamento automatico dei dati per sicurezza multi-tenant.

\subsection{Serializzazione}

Il modello nasconde il campo \texttt{user} nella serializzazione JSON per sicurezza:

\begin{lstlisting}[language=JavaScript]
userActivitySchema.set('toJSON', {
    virtuals: true,
    versionKey: false,
    transform: (doc, ret) => {
        ret.id = ret._id;
        delete ret._id;
        delete ret.user;  // Sicurezza: non esporre userId
        return ret;
    },
});
\end{lstlisting}

\textbf{Punto di Forza}: Prevenzione di data leakage accidentale.

\section{ActivityService}

\subsection{Panoramica}

\texttt{ActivityService} è il servizio centrale che gestisce tutta la logica di registrazione attività e gamification. È implementato come singleton per garantire coerenza.

\subsection{Metodo Principale: recordActivity}

Il metodo \texttt{recordActivity} è il punto di ingresso principale per la registrazione di attività:

\begin{lstlisting}[language=JavaScript, caption=Metodo recordActivity (semplificato)]
async recordActivity(userId, type, payload = {}) {
    // 1. Validazione
    if (!userId) throw AppError.validation('UserId mancante');
    if (!ACTIVITY_TYPES.includes(type)) throw AppError.validation('Tipo non valido');
    
    // 2. Arricchimento metadata
    const enrichedMetadata = {
        ...metadata,
        dayOfWeek: metadata.dayOfWeek || this._getDayOfWeekLabel(createdAt),
        timeOfDay: metadata.timeOfDay || this._getTimeOfDayLabel(createdAt),
    };
    
    // 3. Calcolo XP
    const xpEarned = Math.max(0, this._calculateXp(type, enrichedMetadata));
    
    // 4. Creazione record
    await UserActivity.create({
        user: userId,
        type,
        entityId,
        category,
        sentiment,
        xpEarned,
        metadata: enrichedMetadata,
        createdAt,
    });
    
    // 5. Aggiornamento utente (XP, level, streak, stats)
    // ... (vedi sezione Gamification)
    
    return { leveledUp, xpEarned, newLevel };
}
\end{lstlisting}

\subsection{Arricchimento Automatico Metadata}

Il servizio arricchisce automaticamente i metadata con informazioni contestuali:

\begin{itemize}
    \item \texttt{dayOfWeek}: Etichetta del giorno della settimana (sun, mon, tue, etc.)
    \item \texttt{timeOfDay}: Fascia oraria (morning, afternoon, evening, night)
\end{itemize}

\textbf{Punto di Forza}: Arricchimento automatico facilita analisi comportamentali senza richiedere dati aggiuntivi dai servizi chiamanti.

\subsection{Calcolo XP}

Il sistema calcola gli XP guadagnati in base al tipo di attività:

\begin{lstlisting}[language=JavaScript, caption=Calcolo XP per tipo attività]
_calculateXp(type, metadata) {
    if (type === 'SESSION_COMPLETE') {
        return this._calculateSessionXp(metadata);
    }
    if (type === 'GOAL_COMPLETED') {
        return this._toNumber(metadata.xpEarned, 25);
    }
    if (type === 'STREAK_BONUS') {
        return this._toNumber(metadata.xpEarned, 5);
    }
    return 0;
}
\end{lstlisting}

\subsubsection{Calcolo XP Sessioni di Studio}

Per le sessioni di studio, il calcolo è più complesso:

\begin{lstlisting}[language=JavaScript, caption=Calcolo XP sessioni]
_calculateSessionXp(metadata) {
    const baseXP = 10;
    const correctCount = this._toNumber(metadata.correctCount, 0);
    const timeSpentSeconds = this._toNumber(metadata.timeSpentSeconds, 0);
    const totalCards = this._toNumber(metadata.totalCards, correctCount + wrongCount);
    
    const timePerCard = totalCards > 0 ? timeSpentSeconds / totalCards : 0;
    const timeBonus = timePerCard > 0
        ? Math.max(0, Math.round(10 - timePerCard / 3))
        : 0;
    
    const streakBonus = this._toNumber(metadata.streakBonus, 0);
    
    return baseXP + (correctCount * 2) + timeBonus + streakBonus;
}
\end{lstlisting}

\textbf{Formula}: 
\[
XP = 10 + (2 \times \text{correctCount}) + \max(0, 10 - \frac{\text{timeSpentSeconds}}{3 \times \text{totalCards}}) + \text{streakBonus}
\]

\textbf{Punto di Forza}: Sistema di ricompensa bilanciato che premia accuratezza, velocità e costanza.

\subsection{Gestione Streak Atomica}

\subsubsection{Problema Risolto: Race Condition}

Il sistema utilizza aggiornamenti atomici MongoDB per prevenire race conditions nello streak:

\begin{lstlisting}[language=JavaScript, caption=Update Streak Atomico]
async updateStreakAtomically(userId) {
    const today = this._startOfDay(new Date());
    const yesterdayStart = new Date(today.getTime() - MS_PER_DAY);
    
    // OPERAZIONE 1: Incrementa se lastActivityDate è ieri
    const incrementResult = await User.updateOne(
        {
            _id: userId,
            $or: [
                { 'gamification.streak.lastActivityDate': { $exists: false } },
                { 'gamification.streak.lastActivityDate': null },
                {
                    $and: [
                        { 'gamification.streak.lastActivityDate': { $gte: yesterdayStart } },
                        { 'gamification.streak.lastActivityDate': { $lt: today } },
                    ],
                },
            ],
        },
        {
            $inc: { 'gamification.streak.current': 1 },
            $set: { 'gamification.streak.lastActivityDate': today },
        }
    );
    
    // OPERAZIONE 2: Reset se lastActivityDate è più vecchio di ieri
    if (incrementResult.matchedCount === 0) {
        await User.updateOne(
            {
                _id: userId,
                'gamification.streak.lastActivityDate': {
                    $exists: true,
                    $ne: null,
                    $lt: yesterdayStart,
                },
            },
            {
                $set: {
                    'gamification.streak.current': 1,
                    'gamification.streak.lastActivityDate': today,
                },
            }
        );
    }
    
    // Aggiorna best streak se necessario
    // ...
}
\end{lstlisting}

\textbf{Punto di Forza}: Risoluzione elegante del problema di race condition usando operatori atomici MongoDB invece di read-modify-write.

\textbf{Nota}: Esiste ancora un metodo deprecato \texttt{updateStreak(userDoc)} mantenuto per compatibilità, ma non dovrebbe essere usato.

\subsection{Calcolo Livello}

Il livello viene calcolato dalla formula:

\[
\text{level} = \max(1, \lfloor \sqrt{\frac{\text{xp}}{100}} \rfloor + 1)
\]

\begin{lstlisting}[language=JavaScript]
getLevelFromXP(xp) {
    const safeXp = this._toNumber(xp, 0);
    return Math.max(1, Math.floor(Math.sqrt(safeXp / XP_LEVEL_BASE)) + 1);
}
\end{lstlisting}

\textbf{Punto di Forza}: Formula matematica che garantisce progressione bilanciata (più XP necessari per livelli più alti).

\section{ActivitySubscriber: Pattern Observer}

\subsection{Architettura Decoupling}

Il sistema utilizza il pattern Observer per decoupling tra servizi. I servizi di business emettono eventi invece di chiamare direttamente \texttt{ActivityService}:

\begin{lstlisting}[language=JavaScript, caption=Esempio: WorkLogService emette evento]
// In workLogService.js
async create(tenantScope, data) {
    const created = await super.create(tenantScope, data);
    
    // Pattern Observer: emetti evento invece di chiamare direttamente
    eventBus.emit('worklog.created', {
        userId: this._getUserId(tenantScope),
        workLog: created,
    });
    
    return created;
}
\end{lstlisting}

\textbf{Punto di Forza}: Decoupling completo - i servizi di business non conoscono ActivityService.

\subsection{Eventi Supportati}

L'\texttt{ActivitySubscriber} ascolta i seguenti eventi:

\begin{itemize}
    \item \texttt{worklog.created}: Quando viene creato un nuovo worklog
    \item \texttt{worklog.updated}: Quando viene aggiornato un worklog (solo se cambiamenti significativi)
    \item \texttt{goal.completed}: Quando un goal viene completato
    \item \texttt{session.completed}: Quando una sessione di studio viene completata
\end{itemize}

\subsection{Handler per WorkLog}

\begin{lstlisting}[language=JavaScript, caption=Handler worklog.created]
async function handleWorkLogCreated(data) {
    const { userId, workLog } = data;
    
    // Determina tipo attività
    const activityType = (workLog.startTime && workLog.endTime) 
        ? 'WORK_SESSION_LOGGED' 
        : 'WORK_NOTE_CREATED';
    
    // Usa queue con retry automatico
    await addActivityJob(userId, activityType, {
        entityId: workLog._id || workLog.id,
        category: 'work',
        metadata: {
            durationMinutes: workLog.durationMinutes,
            tags: workLog.tags,
            timeOfDay: getTimeOfDayLabel(new Date()),
            // ...
        },
    });
}
\end{lstlisting}

\textbf{Punto di Forza}: Distinzione intelligente tra sessioni con tracking temporale e note semplici.

\subsection{Inizializzazione}

I subscriber vengono inizializzati all'avvio dell'applicazione:

\begin{lstlisting}[language=JavaScript]
function initializeSubscribers() {
    eventBus.on('worklog.created', handleWorkLogCreated);
    eventBus.on('worklog.updated', handleWorkLogUpdated);
    eventBus.on('goal.completed', handleGoalCompleted);
    eventBus.on('session.completed', handleSessionCompleted);
    
    console.log('✅ ActivitySubscriber inizializzato');
}
\end{lstlisting}

\textbf{⚠️ Problema Potenziale}: Se \texttt{initializeSubscribers()} non viene chiamata all'avvio, gli eventi non verranno processati. Verificare che sia chiamata in \texttt{server/index.js}.

\section{ActivityQueue: Retry Automatico}

\subsection{Architettura Bull Queue}

Il sistema utilizza Bull (basato su Redis) per gestire una coda asincrona con retry automatico:

\begin{lstlisting}[language=JavaScript, caption=Configurazione Queue]
activityQueue = new Queue('activity-tracking', {
    redis: process.env.REDIS_URL || {
        host: process.env.REDIS_HOST || 'localhost',
        port: process.env.REDIS_PORT || 6379,
    },
    defaultJobOptions: {
        attempts: 3,  // 3 tentativi totali
        backoff: {
            type: 'exponential',  // 2s, 4s, 8s
            delay: 2000,
        },
        removeOnComplete: {
            age: 3600,  // Rimuovi dopo 1 ora
            count: 1000,  // Max 1000 job completati
        },
    },
});
\end{lstlisting}

\textbf{Punto di Forza}: Retry automatico con exponential backoff previene perdita di dati in caso di errori temporanei.

\subsection{Fallback a Chiamata Diretta}

Se Redis non è disponibile, il sistema fallback a chiamata diretta:

\begin{lstlisting}[language=JavaScript]
async function addActivityJob(userId, activityType, payload) {
    if (!activityQueue) {
        // Fallback: chiamata diretta (senza retry)
        try {
            await activityService.recordActivity(userId, activityType, payload);
        } catch (error) {
            console.error(`❌ Activity error (direct): ${error.message}`);
            // Non propagare errore (non bloccare l'app)
        }
        return;
    }
    
    // Usa queue con retry
    await activityQueue.add({ userId, activityType, payload });
}
\end{lstlisting}

\textbf{Punto di Forza}: Resilienza - il sistema continua a funzionare anche senza Redis.

\textbf{⚠️ Limitazione}: Senza Redis, non c'è retry automatico. Gli errori vengono solo loggati.

\subsection{Monitoring}

La queue emette eventi per monitoring:

\begin{lstlisting}[language=JavaScript]
activityQueue.on('completed', (job, result) => {
    eventBus.emit('activity.queue.completed', {
        jobId: job.id,
        activityType: job.data.activityType,
        attempts: job.attemptsMade,
    });
});

activityQueue.on('failed', (job, error) => {
    eventBus.emit('activity.queue.failed', {
        jobId: job.id,
        activityType: job.data.activityType,
        attempts: job.attemptsMade,
        error: error.message,
    });
});
\end{lstlisting}

\textbf{Punto di Forza}: Tracciabilità completa per debugging e monitoring.

\section{Integrazioni con Altri Servizi}

\subsection{Servizi che Registrano Attività}

I seguenti servizi registrano attività direttamente (non tramite eventi):

\begin{enumerate}
    \item \textbf{GoalService}: Registra \texttt{GOAL_CREATED} quando viene creato un obiettivo
    \item \textbf{StudyService}: Registra \texttt{SESSION_COMPLETE} quando una sessione di studio viene completata
    \item \textbf{ProjectService}: Registra \texttt{PROJECT_CREATED} e \texttt{PROJECT_COMPLETED}
    \item \textbf{CheckInService}: Registra \texttt{HABIT_CHECKIN} quando viene fatto un check-in
\end{enumerate}

\textbf{⚠️ Inconsistenza Architetturale}: Alcuni servizi chiamano direttamente \texttt{activityService.recordActivity()}, altri usano il pattern Observer. Questo crea inconsistenza:

\begin{itemize}
    \item \texttt{WorkLogService} usa eventi (corretto)
    \item \texttt{GoalService} chiama direttamente (inconsistente)
    \item \texttt{StudyService} chiama direttamente (inconsistente)
\end{itemize}

\textbf{Raccomandazione}: Standardizzare su pattern Observer per tutti i servizi.

\section{Utilizzo in Dashboard e Analytics}

\subsection{Dashboard Controller}

Il \texttt{dashboardController} utilizza \texttt{UserActivity} per:

\begin{itemize}
    \item \textbf{Jump Back In}: Ultime 10 attività per suggerire continuazione
    \item \textbf{Statistiche}: Conteggio attività per categoria
\end{itemize}

\begin{lstlisting}[language=JavaScript, caption=Query attività recenti per dashboard]
const recentActivities = await UserActivity.find({ user: userId })
    .sort({ createdAt: -1 })
    .limit(10)
    .lean();
\end{lstlisting}

\textbf{Punto di Forza}: Query efficiente grazie agli indici.

\subsection{Analytics Controller}

Il \texttt{analyticsController} aggrega dati per visualizzazioni settimanali:

\begin{lstlisting}[language=JavaScript, caption=Aggregazione settimanale]
const dailyRaw = await UserActivity.aggregate([
    {
        $match: {
            user: userId,
            createdAt: { $gte: startDate },
        },
    },
    {
        $group: {
            _id: { $dateToString: { format: '%Y-%m-%d', date: '$createdAt' } },
            totalXP: { $sum: '$xpEarned' },
            studySeconds: { $sum: { /* ... */ } },
            workMinutes: { $sum: { /* ... */ } },
            goalsCompleted: { $sum: { /* ... */ } },
        },
    },
    { $sort: { _id: 1 } },
]);
\end{lstlisting}

\textbf{Punto di Forza}: Aggregazione efficiente con pipeline MongoDB.

\subsection{Insight Service}

Il \texttt{insightService} analizza gli ultimi 30 giorni di attività per generare consigli:

\begin{lstlisting}[language=JavaScript, caption=Query per insights]
const activities = await UserActivity.forTenant(userId)
    .find({ createdAt: { $gte: startDate } })
    .select('type metadata createdAt')
    .lean();
\end{lstlisting}

\textbf{Analisi Eseguite}:

\begin{enumerate}
    \item \textbf{Cronotipo}: Analizza precisione per fascia oraria (morning vs night)
    \item \textbf{Ambition Gap}: Confronta obiettivi creati vs completati
    \item \textbf{Consistency Streak}: Identifica giorni della settimana senza attività
\end{enumerate}

\textbf{Punto di Forza}: Analisi comportamentale avanzata basata su dati reali.

\section{Punti di Forza}

\subsection{Architettura}

\begin{enumerate}
    \item \textbf{Decoupling}: Pattern Observer separa logica di business da tracking
    \item \textbf{Resilienza}: Fallback a chiamata diretta se Redis non disponibile
    \item \textbf{Atomicità}: Aggiornamenti streak atomici prevengono race conditions
    \item \textbf{Multi-Tenancy}: Isolamento automatico dei dati per sicurezza
\end{enumerate}

\subsection{Performance}

\begin{enumerate}
    \item \textbf{Indici Ottimizzati}: Query efficienti per pattern comuni
    \item \textbf{Asincronia}: Queue non blocca risposte HTTP
    \item \textbf{Aggregazioni}: Pipeline MongoDB efficienti per analytics
\end{enumerate}

\subsection{Sicurezza}

\begin{enumerate}
    \item \textbf{Isolamento Dati}: Plugin multi-tenancy previene data leakage
    \item \textbf{Serializzazione Sicura}: Campo \texttt{user} nascosto in JSON
    \item \textbf{Validazione}: Controlli su tipo attività e payload
\end{enumerate}

\subsection{Manutenibilità}

\begin{enumerate}
    \item \textbf{Codice Pulito}: Separazione responsabilità chiara
    \item \textbf{Documentazione}: Commenti esplicativi nel codice
    \item \textbf{Testabilità}: Architettura facilita unit testing
\end{enumerate}

\section{Problemi e Limitazioni}

\subsection{Bug Conosciuti}

\subsubsection{Bug 1: Inconsistenza Pattern Observer}

\textbf{Problema}: Alcuni servizi chiamano direttamente \texttt{activityService.recordActivity()}, altri usano eventi.

\textbf{Impatto}: 
\begin{itemize}
    \item Inconsistenza architetturale
    \item Difficile tracciare tutte le sorgenti di attività
    \item Testing più complesso
\end{itemize}

\textbf{Severità}: Media

\textbf{Soluzione Proposta}: Refactoring per standardizzare su pattern Observer.

\subsubsection{Bug 2: DAY\_LABELS Incompleto}

Nel file \texttt{activityService.js}, la costante \texttt{DAY\_LABELS} è definita come:

\begin{lstlisting}[language=JavaScript]
const DAY_LABELS = ['sun', 'mon', 'tue', 'wed', 'thu', 'fri', 'sat'];
\end{lstlisting}

Ma in alcuni punti del codice (riga 29) appare:

\begin{lstlisting}[language=JavaScript]
const DAY_LABELS = [ 'fri', 'sat'];
\end{lstlisting}

\textbf{Problema}: Array incompleto causerebbe errori in \texttt{\_getDayOfWeekLabel()}.

\textbf{Impatto}: 
\begin{itemize}
    \item Errori in produzione per giorni mancanti
    \item Metadata incompleti
\end{itemize}

\textbf{Severità}: Alta (se presente nel codice di produzione)

\textbf{Soluzione}: Verificare e correggere la definizione.

\subsubsection{Bug 3: Metodo Deprecato Ancora Presente}

Il metodo \texttt{updateStreak(userDoc)} è marcato come \texttt{@deprecated} ma ancora presente:

\begin{lstlisting}[language=JavaScript]
/**
 * @deprecated Usa updateStreakAtomically() invece
 * Mantenuto per compatibilità temporanea
 */
updateStreak(userDoc) {
    // ... implementazione non atomica
}
\end{lstlisting}

\textbf{Problema}: Metodo non atomico ancora disponibile potrebbe essere usato per errore.

\textbf{Impatto}: 
\begin{itemize}
    \item Possibili race conditions se usato
    \item Confusione per sviluppatori
\end{itemize}

\textbf{Severità}: Bassa (se non usato)

\textbf{Soluzione}: Rimuovere dopo verifica che non sia usato.

\subsection{Limitazioni}

\subsubsection{Limitazione 1: Nessun Retry Senza Redis}

\textbf{Problema}: Se Redis non è disponibile, il fallback a chiamata diretta non ha retry.

\textbf{Impatto}: 
\begin{itemize}
    \item Perdita di attività in caso di errori temporanei
    \item Nessuna garanzia di delivery
\end{itemize}

\textbf{Severità}: Media

\textbf{Soluzione Proposta}: Implementare retry locale (in-memory queue) come fallback.

\subsubsection{Limitazione 2: Nessuna Validazione Metadata}

\textbf{Problema}: Il campo \texttt{metadata} è di tipo \texttt{Mixed} senza validazione schema.

\textbf{Impatto}: 
\begin{itemize}
    \item Possibili dati inconsistenti
    \item Difficile debugging
    \item Nessuna garanzia di struttura
\end{itemize}

\textbf{Severità}: Bassa-Media

\textbf{Soluzione Proposta}: Validazione runtime o schema Mongoose per metadata comuni.

\subsubsection{Limitazione 3: Nessuna Archiviazione}

\textbf{Problema}: Le attività vengono mantenute indefinitamente nel database.

\textbf{Impatto}: 
\begin{itemize}
    \item Crescita database nel tempo
    \item Query più lente su dataset grandi
    \item Costi storage crescenti
\end{itemize}

\textbf{Severità}: Bassa (a breve termine), Alta (a lungo termine)

\textbf{Soluzione Proposta}: 
\begin{itemize}
    \item Archiviazione attività vecchie (> 1 anno) in collection separata
    \item Rimozione attività molto vecchie (> 3 anni) o export a data warehouse
\end{itemize}

\subsubsection{Limitazione 4: Nessun Rate Limiting}

\textbf{Problema}: Non c'è protezione contro spam di attività (es. loop infiniti).

\textbf{Impatto}: 
\begin{itemize}
    \item Possibile DoS involontario
    \item Database growth anomalo
    \item Performance degradation
\end{itemize}

\textbf{Severità}: Media

\textbf{Soluzione Proposta}: Rate limiting per tipo attività per utente (es. max 100 attività/minuto).

\subsubsection{Limitazione 5: Timezone Non Gestito}

\textbf{Problema}: \texttt{dayOfWeek} e \texttt{timeOfDay} sono calcolati sul server, non sul timezone utente.

\textbf{Impatto}: 
\begin{itemize}
    \item Analisi comportamentali inaccurate per utenti in timezone diversi
    \item Insights potenzialmente errati
\end{itemize}

\textbf{Severità}: Media

\textbf{Soluzione Proposta}: Salvare timezone utente e calcolare \texttt{dayOfWeek}/\texttt{timeOfDay} in base a quello.

\subsection{Errori Potenziali}

\subsubsection{Errore 1: Race Condition in Statistiche}

Nel metodo \texttt{recordActivity}, le statistiche vengono aggiornate con read-modify-write:

\begin{lstlisting}[language=JavaScript]
const rawStats = gamification.stats || {};
const nextStats = {
    totalStudySessions: this._toNumber(rawStats.totalStudySessions, 0),
    // ...
};

if (type === 'SESSION_COMPLETE') {
    nextStats.totalStudySessions += 1;
    // ...
}

user.set('gamification', {
    // ...
    stats: nextStats,
});
await user.save();
\end{lstlisting}

\textbf{Problema}: Non atomico - possibile perdita di aggiornamenti in caso di richieste simultanee.

\textbf{Impatto}: Statistiche inaccurate

\textbf{Severità}: Media

\textbf{Soluzione Proposta}: Usare operatori atomici MongoDB (\texttt{\$inc}) per statistiche.

\subsubsection{Errore 2: Nessun Rollback su Errore}

Se \texttt{UserActivity.create()} riesce ma \texttt{user.save()} fallisce, l'attività rimane registrata ma l'XP non viene aggiornato.

\textbf{Problema}: Inconsistenza tra UserActivity e User.gamification

\textbf{Impatto}: 
\begin{itemize}
    \item XP registrati in UserActivity ma non in User
    \item Statistiche inconsistenti
\end{itemize}

\textbf{Severità}: Media

\textbf{Soluzione Proposta}: Transazione MongoDB o compensazione (rollback UserActivity se save fallisce).

\section{Raccomandazioni per Miglioramenti}

\subsection{Breve Termine}

\begin{enumerate}
    \item \textbf{Standardizzare Pattern Observer}: Refactoring per usare eventi ovunque
    \item \textbf{Correggere DAY\_LABELS}: Verificare e correggere definizione
    \item \textbf{Rimuovere Metodo Deprecato}: Eliminare \texttt{updateStreak(userDoc)}
    \item \textbf{Aggiungere Rate Limiting}: Protezione contro spam
    \item \textbf{Validazione Metadata}: Schema per metadata comuni
\end{enumerate}

\subsection{Medio Termine}

\begin{enumerate}
    \item \textbf{Retry Locale}: Implementare in-memory queue come fallback
    \item \textbf{Statistiche Atomiche}: Usare operatori MongoDB per aggiornamenti
    \item \textbf{Transazioni}: Implementare transazioni per consistenza
    \item \textbf{Timezone Support}: Calcolo dayOfWeek/timeOfDay basato su timezone utente
    \item \textbf{Monitoring Avanzato}: Dashboard per metriche queue e errori
\end{enumerate}

\subsection{Lungo Termine}

\begin{enumerate}
    \item \textbf{Archiviazione}: Sistema di archiviazione per attività vecchie
    \item \textbf{Data Warehouse}: Export periodico per analisi avanzate
    \item \textbf{Streaming}: Real-time analytics con Kafka/Redis Streams
    \item \textbf{ML Integration}: Machine learning per predizioni comportamentali
    \item \textbf{A/B Testing}: Framework per testare algoritmi gamification
\end{enumerate}

\section{Conclusioni}

Il sistema \texttt{UserActivity} è un componente ben architettato e robusto che fornisce funzionalità essenziali per gamification e analytics. I punti di forza principali sono:

\begin{itemize}
    \item Architettura decoupled con pattern Observer
    \item Resilienza con fallback e retry automatico
    \item Atomicità per operazioni critiche (streak)
    \item Performance ottimizzate con indici e aggregazioni
\end{itemize}

Tuttavia, esistono alcune aree di miglioramento:

\begin{itemize}
    \item Inconsistenza nell'uso del pattern Observer
    \item Mancanza di retry senza Redis
    \item Possibili race conditions in statistiche
    \item Nessuna archiviazione per dati storici
\end{itemize}

Con le raccomandazioni proposte, il sistema può essere ulteriormente migliorato per garantire scalabilità, affidabilità e manutenibilità a lungo termine.

\section{Appendice: Tipi di Attività}

\begin{longtable}{|p{4cm}|p{3cm}|p{8cm}|}
\hline
\textbf{Tipo} & \textbf{XP} & \textbf{Descrizione} \\
\hline
\endfirsthead
\hline
\textbf{Tipo} & \textbf{XP} & \textbf{Descrizione} \\
\hline
\endhead
\hline
\endfoot
\hline
\endlastfoot
\texttt{SESSION\_COMPLETE} & Variabile & Sessione di studio completata. XP calcolati in base a correttezza, velocità e streak bonus. \\
\hline
\texttt{GOAL\_CREATED} & 0 & Obiettivo creato. Non genera XP ma aggiorna streak. \\
\hline
\texttt{GOAL\_COMPLETED} & 25 (default) & Obiettivo completato. XP configurabili nel metadata. \\
\hline
\texttt{GOAL\_ARCHIVED} & 0 & Obiettivo archiviato. Non genera XP. \\
\hline
\texttt{WORK\_SESSION\_LOGGED} & 0 & Sessione di lavoro registrata con tracking temporale. Non genera XP ma aggiorna streak. \\
\hline
\texttt{WORK\_NOTE\_CREATED} & 0 & Nota di lavoro creata (senza tracking temporale). Non genera XP ma aggiorna streak. \\
\hline
\texttt{HABIT\_CHECKIN} & 0 & Check-in su abitudine. Non genera XP ma aggiorna streak. \\
\hline
\texttt{HABIT\_SKIPPED} & 0 & Abitudine saltata. Non genera XP. \\
\hline
\texttt{PROJECT\_CREATED} & 0 & Progetto creato. Non genera XP ma aggiorna streak. \\
\hline
\texttt{PROJECT\_COMPLETED} & 0 & Progetto completato. Non genera XP ma aggiorna streak. \\
\hline
\texttt{STREAK\_BONUS} & 5 (default) & Bonus per mantenimento streak. XP configurabili nel metadata. \\
\hline
\end{longtable}

\section{Appendice: Esempi di Metadata}

\subsection{SESSION\_COMPLETE}

\begin{lstlisting}[language=JSON]
{
    "correctCount": 15,
    "wrongCount": 3,
    "totalCards": 18,
    "timeSpentSeconds": 450,
    "streakBonus": 5,
    "deckId": "507f1f77bcf86cd799439011",
    "mode": "review",
    "dayOfWeek": "mon",
    "timeOfDay": "afternoon"
}
\end{lstlisting}

\subsection{GOAL\_COMPLETED}

\begin{lstlisting}[language=JSON]
{
    "entityId": "507f1f77bcf86cd799439012",
    "category": "fitness",
    "priority": "high",
    "daysToComplete": 7,
    "deadline": "2024-01-15T00:00:00.000Z",
    "timeOfDay": "evening",
    "xpEarned": 50
}
\end{lstlisting}

\subsection{WORK\_SESSION\_LOGGED}

\begin{lstlisting}[language=JSON]
{
    "entityId": "507f1f77bcf86cd799439013",
    "projectId": "507f1f77bcf86cd799439014",
    "durationMinutes": 120,
    "tags": ["development", "backend"],
    "timeOfDay": "morning",
    "date": "2024-01-15",
    "startTime": "09:00",
    "endTime": "11:00",
    "hasTimeTracking": true
}
\end{lstlisting}

\end{document}
